\section{Conclusions, limitations and plans}
\label{sec:conclusions}

\NHFOAM\ follows NHWAVE's pressure arrangement on OpenFOAM's unstructured,
moving $\sigma$-meshes. The tests of \citet{ma2012} were run with both models
on the same set-ups, from the simplest upwards, and they changed the code:
\begin{itemize}[nosep]
  \item \textbf{Time accuracy.} With one outer corrector the convection term
        was linearised with the flux of the previous time level, which made the
        scheme first order in time for nonlinear waves. The convecting flux is
        now extrapolated from the last two levels. A solitary wave of
        $H/h=0.5$ lost 6.4\% of its height in 150 m before and 1.0\% after the
        fix (NHWAVE 2.3\%), at no extra cost.
  \item \textbf{Convection.} The TVD scheme damps steep short waves even when
        they do not break; filtered central differencing
        (\code{filteredLinear2V 0.2 0}) is now recommended for non-breaking
        waves. It reduced the standing-wave error to 0.16--0.25\% for
        $kD\ge 2\pi$ (NHWAVE 0.6--4.8\%) and stayed stable in 40--65 s runs over
        the submerged bar. Breaking still needs the TVD scheme.
  \item \textbf{Consistency with NHWAVE.} After the fix, the phase of the waves
        over and behind the submerged bar agrees with NHWAVE's. The closer
        match of the measured phase before the fix came from the first-order
        lag, which slowed the waves by about 0.5\%.
  \item \textbf{Resolution.} On the elliptic shoal, with twice the paper's
        grid spacing and three layers, \NHFOAM's errors at the first sections
        (0.06 and 0.13) are close to NHWAVE's on the paper's grid (0.07 and
        0.11), while NHWAVE on the coarse grid loses 30\% of the incident wave
        height before the shoal.
  \item \textbf{Robustness.} A domain in which waves break anywhere needs the
        TVD scheme and the breaking model, even outside the region of
        interest: the shoal run with the filtered scheme failed when the waves
        broke in a shallow corner, and ran on with \code{limitedLinearV} and
        the hydrostatic front approximation.
\end{itemize}

\paragraph{Limitations.}
\begin{itemize}[nosep]
  \item \textbf{No wetting and drying.} Run-up, inundation and the setup at the
        shoreline cannot be modelled yet; beaches must end in a wall or a
        shelf with a sponge. Tests 4.5 and 4.7 of \citet{ma2012} wait for it.
  \item \textbf{Breaking} has been tested on one case. The onset is right, but
        the hydrostatic bores lose too little energy in the surf zone, where
        the wave heights are up to twice the measured ones. Breaking columns
        spread only to neighbours on the same processor.
  \item \textbf{Wave generation} uses linear targets in relaxation zones only.
        Steep or long waves carry spurious free harmonics. There is no
        boundary wavemaker and no irregular, cnoidal or stream-function input.
  \item \textbf{Long waves in closed flumes.} The zones absorb long waves only
        partly and the wave-making zone does not conserve volume, so long waves
        of a few millimetres build up and shift the phase over shallow bars in
        long runs.
  \item \textbf{Restarts} do not save the time history of $\eta$ and
        $\bm{u}_f$; the first step after a restart is Euler and first order.
  \item \textbf{Landslide trailing waves} are 1.5--3 times the measured ones
        (Sections~\ref{sec:seabed} and \ref{sec:ma-landslide}), probably
        because a wall replaces the beach.
\end{itemize}

\paragraph{Plans.} In order of priority:
\begin{enumerate}[nosep]
  \item wetting and drying with thin-film columns and $q$ switched off in very
        shallow water, then tests 4.5 and 4.7 of \citet{ma2012};
  \item breaking: the dissipation of the bores (momentum-conserving face depths
        at the fronts, or a small eddy viscosity there), a kinematic onset
        criterion and stabilised $k$--$\omega$ \citep{larsen2018} for
        many-layer runs;
  \item volume-conserving wave-making zones and absorption of long waves;
  \item nonlinear and irregular wave generation and a boundary wavemaker;
  \item depth-integrated granular slides for landslide tsunamis;
  \item coupling to \code{interFoam} and \code{sedInterFoam} through relaxation
        zones, and hybrid meshes around structures.
\end{enumerate}
